\color{unidarkgreen}{\subsection{Фиксация положения выключателей}}
\color{black}

\begin{enumerate}[label=\arabic{section}.\arabic{subsection}.\arabic*,  align=left, labelsep=4pt, leftmargin=0pt, itemindent=57pt]

\item
Функциональный блок фиксации положения выключателей (ФПВ) предназначен для определения
состояния выключателей на основании технологической информации: 
\begin{itemize}
\item состояния блок-контактов выключателя;
\item положения разъединителей в цепи выключателя;
\item информации о выводе в ремонт выключателя.
\end{itemize}

\item
Контроль положения разъединителей при формировании сигналов включенного или отключенного
положения выключателей вводится программным переключателем \textbf{«Контр Р В1»}.  Также имеется
возможность выбрать количество разъединителей, участвующих в формировании сигналов включенного или
отключенного положения выключателей, задается программным переключателем \textbf{«Схема Р В1»}. А для
выключателей В2 своей и смежной стороны имеется возможность выбрать тип выключателя линейный или
обходной, задается программным переключателем \textbf{«Тип В2»}.

\item
Также функция ФПВ пофазно контролирует состояние блок-контактов выключателя и в случае
расхождения полюсов формирует сигнал \textit{«Внутр. пуск. ЗНФ В1»}.

\item
Ввод ФПВ остальных выключателей и настройка параметров срабатывания осуществляется аналогично.

\begin{figure}[H]
\centering
\includegraphics[width=0.8\textwidth,height=0.8\textheight,keepaspectratio]{\fbpath/761.3381 ФПВ/img/FPV.pdf}
\captionof{figure}{Функциональная схема ФПВ}\label{pic_fpv:FPV}
\end{figure}

\small
\begin{longtable}{|>{\centering\arraybackslash}m{5.3cm}|>{\centering\arraybackslash}m{4.2cm}|>{\centering\arraybackslash}m{1cm}|>{\centering\arraybackslash}m{3.3cm}|>{\centering\arraybackslash}m{1.8cm}|}
% Первая страница: полный заголовок
\caption{Перечень уставочных параметров ФПВ\hfill\vspace{-0.5\baselineskip}}\label{fpv:tbl1}\\
\hline
\rowcolor{gray!30}
Наименование уставки & Диапазон & Шаг & Значение по умолчанию & Единица измерения \\
\hline
\endfirsthead

% Промежуточные страницы: «Продолжение таблицы»
\caption*{\hspace{3pt}\emph{Продолжение таблицы \ref{fpv:tbl1}}\hfill\ }\vspace{-0.5\baselineskip} \\
\hline
\rowcolor{gray!30}
Наименование уставки & Диапазон & Шаг & Значение по умолчанию & Единица измерения \\
\hline
\endhead

%===m>ФПВ|-
%===m
\end{longtable}
\normalsize	

\end{enumerate}